La estructura del documento a partir de aquí es la siguiente:

En el Capítulo \ref{CAP:EST_ARTE} se presenta el estado del arte. Este Cap, donde se introducen los conceptos necesarios para contextualizar el trabajo, incluyendo los sistemas de detección de intrusiones, los conjuntos de datos utilizados en este ámbito y las técnicas de aprendizaje automático aplicadas a la detección de ataques en red.

En el Capítulo 3 se describe la metodología seguida durante el desarrollo del trabajo. En este capítulo se explica la selección del conjunto de datos, la definición de nuevas características, el proceso de extracción y enriquecimiento de los datos, así como el entrenamiento y evaluación del modelo de clasificación.

En el Capítulo 4 se muestran los resultados obtenidos. Primero se analiza el conjunto de datos final y, posteriormente, se presentan los experimentos realizados, junto con sus métricas de evaluación, matrices de confusión e importancia de características. Además, se comparan los resultados obtenidos en ambos experimentos.
